\subsection{相对论、引力与宇宙学}

\begin{结论}{Lorentz 变换的三个推论}
$t'=\gamma\bigl(t-\frac{vx}{c^2}\bigr)$，$x'=\gamma(x-vt)$；间隔 $s^2=c^2\Delta t^2-\Delta x^2$ 不变。
\[
\text{时间延缓 }\Delta\tau=\frac{\Delta t}{\gamma},\quad \text{长度收缩 }L=\frac{L_0}{\gamma},\quad \重点{\text{同时的相对性：}\Delta t'=-\gamma\frac{v\Delta x}{c^2}}.
\]
$\mu$ 子 $\tau=2.2\,\mu$s，$E=3\,$GeV（$\gamma\approx28$）时衰变长度 $\gamma\beta c\tau\approx19\,$km，故能从高空到达地面。多数“佯谬”（车库与长杆、双生子）都靠同时性相对性化解。
\end{结论}

\begin{结论}{质心系速度与动量变换}
多粒子系统的零动量系速度：
\[
\重点{\boldsymbol\beta_{\mathrm{cm}}=\frac{\sum\mathbf p_ic}{\sum E_i}},\qquad E'=\gamma(E-\boldsymbol\beta\cdot\mathbf pc),\quad p'_\parallel=\gamma\Bigl(p_\parallel-\frac{\beta E}{c}\Bigr),\quad p'_\perp=p_\perp .
\]
$(E/c,\mathbf p)$ 与 $(ct,\mathbf x)$ 同样变换。光子：$E'=E\gamma(1-\beta\cos\theta)$，即 Doppler 公式。
\end{结论}

\begin{结论}{相对论弹性碰撞与 Mandelstam 变量}
等质量粒子弹性碰撞（一个静止）：两出射粒子夹角 \重点{\theta_1+\theta_2<90^\circ}（非相对论极限才等于 $90^\circ$），可在气泡室照片中直接判断相对论效应。
二体过程 $s=(p_1+p_2)^2$，$t=(p_1-p_3)^2$，$u=(p_1-p_4)^2$（$c=1$）：$s+t+u=\sum_{i=1}^4m_i^2$；$-t$ 为动量转移的平方，决定散射角分布。
\end{结论}

\begin{结论}{相对论火箭}
喷射速度 $u$（相对火箭），从静止加速到 $\beta$：
\[
\重点{\frac{m_i}{m_f}=\Bigl(\frac{1+\beta}{1-\beta}\Bigr)^{c/2u}}\quad\Leftrightarrow\quad \eta=\frac uc\ln\frac{m_i}{m_f}\ (\text{快度}).
\]
光子火箭（$u=c$）达到 $\beta=0.9$ 需 $m_i/m_f=\sqrt{19}\approx4.4$；$u\ll c$ 时回到 Tsiolkovsky 公式。
\end{结论}

\begin{结论}{由动能求速度}
\[
\重点{\beta=\sqrt{1-\Bigl(\frac{mc^2}{T+mc^2}\Bigr)^2}},\qquad pc=\sqrt{T^2+2Tmc^2}.
\]
$1\,$MeV 电子：$\gamma=2.96$，$\beta=0.941$；$1\,$MeV 质子：$\beta=0.046$。$T\ll mc^2$ 时 $T=\frac12mv^2+\frac38mv^4/c^2$，相对论修正比例 $\frac34\beta^2$。
\end{结论}

\begin{结论}{Schwarzschild 时空的特征半径}
$r_s=\frac{2GM}{c^2}$（太阳 $2.95\,$km）：
\[
\text{光子球 }r=\tfrac32r_s,\qquad \重点{\text{最内稳定圆轨道 ISCO }r=3r_s=\frac{6GM}{c^2}},\qquad \text{ISCO 处结合能 }\approx5.7\%\,mc^2 .
\]
远处观察者看到静止钟变慢 $\sqrt{1-r_s/r}$。吸积盘辐射效率约 $6\%$（自转黑洞可达 $42\%$），远高于核聚变的 $0.7\%$。
\end{结论}

\begin{结论}{Hawking 温度与黑洞寿命}
\[
\重点{T_H=\frac{\hbar c^3}{8\pi GMk}}\approx6\times10^{-8}\,\mathrm K\cdot\frac{M_\odot}{M},\qquad S_{\mathrm{BH}}=\frac{kc^3A}{4G\hbar}.
\]
质量越小越热（负热容）；寿命 $\propto M^3$，太阳质量黑洞约 $10^{67}$ 年。熵与\textbf{面积}而非体积成正比（全息原理）。
\end{结论}

\begin{结论}{Hubble 定律与临界密度}
$v=H_0d$，$H_0\approx70\,$km/s/Mpc：Hubble 时间 $1/H_0\approx14\,$Gyr，Hubble 半径 $c/H_0\approx4.3\,$Gpc。
\[
\重点{\rho_c=\frac{3H_0^2}{8\pi G}\approx9\times10^{-27}\,\mathrm{kg/m^3}}\ (\text{约每立方米 5 个质子}).
\]
红移 $1+z=a_0/a$；辐射温度 $T\propto1/a$（CMB 今天 $2.725\,$K，复合时 $z\approx1100$、$T\approx3000\,$K）。
\end{结论}

\begin{结论}{Friedmann 方程与各成分的标度}
\[
\重点{H^2=\frac{8\pi G}{3}\rho-\frac{kc^2}{a^2}},\qquad \rho_{\mathrm m}\propto a^{-3},\quad \rho_{\mathrm r}\propto a^{-4},\quad \rho_\Lambda=\text{常数}.
\]
平坦、单一成分：物质主导 $a\propto t^{2/3}$，辐射主导 $a\propto t^{1/2}$，$\Lambda$ 主导 $a\propto\ee^{Ht}$。物质主导宇宙年龄 $\frac23H_0^{-1}$，与恒星年龄的矛盾是暗能量的早期证据之一。
\end{结论}

\begin{结论}{引力透镜的 Einstein 半径}
点质量 $M$，观察者--透镜、观察者--源、透镜--源距离 $D_L,D_S,D_{LS}$：
\[
\重点{\theta_E=\sqrt{\frac{4GM}{c^2}\frac{D_{LS}}{D_LD_S}}}.
\]
源恰在透镜后方时成 Einstein 环；偏离时成两个像，总放大率在 $u=\beta/\theta_E=1$ 处为 $1.34$（微引力透镜探测阈值）。银河系内恒星透镜 $\theta_E\sim$ 毫角秒。
\end{结论}

\参考{Taylor--Wheeler, \emph{Spacetime Physics}；Hartle, \emph{Gravity}；Carroll, \emph{Spacetime and Geometry}；Ryden, \emph{Introduction to Cosmology}；Hagedorn, \emph{Relativistic Kinematics}（MIT 8.033 / 8.286 / 8.962）。}
